% Part: incompleteness
% Chapter: theories-computability
% Section: q-is-ce

\documentclass[../../../include/open-logic-section]{subfiles}

\begin{document}

\olfileid{inc}{tcp}{qce}

\olsection{$\Th{Q}$ 是可计算枚举完备的}


\begin{thm}
$\Th{Q}$ 是 c.e. 但不是可判定的。事实上，它是一个完全的 c.e. 集。
\end{thm}

\begin{proof}
不难看出 $\Th{Q}$ 是 c.e.，因为它是所有满足如下条件的语句 $y$（的编码）的集合：存在一个在 $\Th{Q}$ 中关于 $y$ 的证明 $x$：
\[
Q = \Setabs{y}{\lexists[x][\Prf[\Th{Q}](x,y)]}.
\]
但我们知道 $\Prf[\Th{Q}](x,y)$ 是可计算的（事实上是原始递归的），而任何能写成上述形式的集合都是 c.e.。

它是一个完全的 c.e. 集，等价于 $K \leq_m Q$，其中 $K = \Setabs{x}{!A_x(x) \downarrow}$。下面证明 $K$ 可归约到 $\Th{Q}$。由于 Kleene 谓词 $T(e,x,s)$ 是原始递归的，它在 $\Th{Q}$ 中是可表示的，设由 $!A_T$ 表示。那么对任意 $x$，我们有
\begin{align*}
x \in K & \lif \lexists[s][T(x,x,s)] \\
& \lif \lexists[s][(\Th{Q} \Proves !A_T(\num x,\num x, \num s))]
  \\
& \lif \Th{Q} \Proves \lexists[s][!A_T(\num x, \num x, s)].
\end{align*}
反过来，如果 $\Th{Q} \Proves \lexists[s][!A_T(\num x, \num x, s)]$，那么实际上存在某个自然数 $n$ 使得公式 $!A_T(\num x,
\num x, \num n)$ 为真。由于 $!A_T$ 表示 $T$，若 $T(x,x,n)$ 为假，则 $\Th{Q}$ 将证明 $\lnot !A_T(\num x, \num x, \num n)$。但这样 $\Th{Q}$ 就证明了一个假公式，产生矛盾。所以 $T(x,x,n)$ 必然为真，这意味着 $!A_x(x) \downarrow$。

简言之，对任意 $x$，$x \in K$ 当且仅当 $\Th{Q}$ 证明 $\lexists[s][T(\num x,\num x,s)]$。所以，将 $x$ 映射到语句 $\lexists[s][T(\num x,
  \num x, s)]$（的编码）的函数 $f$，是 $K$ 到 $\Th{Q}$ 的一个归约。
\end{proof}

\end{document}
